% ============================================================
% Déclaration d'usage de l'IA proposée (agent 7)
% Remplace intégralement frontmatter/declaration_ia.tex.
% Les champs [À COMPLÉTER PAR L'AUTEUR] doivent être renseignés ou supprimés
% avant compilation ; ne laisser aucune mention entre crochets dans la version
% déposée.
% ============================================================

\chapter*{Déclaration relative à l'utilisation d'outils d'intelligence artificielle}
\addcontentsline{toc}{chapter}{Déclaration relative à l'utilisation d'outils d'intelligence artificielle}

Des outils d'intelligence artificielle générative ont été utilisés comme
assistants au cours de ce travail. La présente déclaration précise les outils,
les parties concernées, la nature de l'aide et la manière dont les
propositions ont été vérifiées.

\section*{Outils utilisés}

\begin{table}[H]
\centering
\small
\begin{tabularx}{\textwidth}{p{0.27\textwidth} p{0.20\textwidth} X}
\toprule
\textbf{Outil (éditeur)} & \textbf{Période} & \textbf{Usage principal} \\
\midrule

elax [À COMPLÉTER PAR L'AUTEUR : modèle GPT et interface utilisés, p.~ex. ChatGPT, version] (OpenAI) &
[À COMPLÉTER] &
Relecture critique simulée du manuscrit, reformulation, aide à la structuration. \\

elax [À COMPLÉTER PAR L'AUTEUR : confirmer l'usage d'OpenAI Codex, suggéré par la configuration \texttt{.codex/} et \texttt{AGENTS.md} du dépôt] &
[À COMPLÉTER] &
Assistance à l'écriture, à la révision et au test de scripts Python et PowerShell. \\
Claude Code, modèle Claude Opus (Anthropic) &
[À COMPLÉTER : dates] 2026 &
Révision finale : audit de cohérence entre texte, code et artefacts, propositions de correction rédactionnelle, préparation des fichiers de publication du dépôt. \\

elax [À COMPLÉTER PAR L'AUTEUR : outil désigné « Luna » dans les notes de travail, ou supprimer cette ligne] &
[À COMPLÉTER] &
[À COMPLÉTER] \\
\bottomrule
\end{tabularx}
\end{table}

\section*{Parties concernées et nature de l'aide}

\begin{itemize}
  \item \textbf{Code du prototype et scripts expérimentaux}
  (\texttt{src/logminer/}, \texttt{scripts/}, \texttt{tests/}) : génération de
  propositions de code, refactorisation, écriture de tests, aide au
  débogage. [À COMPLÉTER PAR L'AUTEUR : préciser les modules principalement
  rédigés par l'auteur et ceux dont une première version a été proposée par un
  outil.]

  \item \textbf{Exécution des campagnes} : les campagnes ont été lancées sur
  les machines de l'auteur ; les artefacts CSV et JSON cités sont produits par
  les scripts du dépôt et non rédigés par un outil. [À COMPLÉTER PAR
  L'AUTEUR : indiquer si un agent logiciel a déclenché lui-même certaines
  exécutions.]

  \item \textbf{Rédaction du manuscrit} : aide à la structuration des
  chapitres, reformulation, harmonisation terminologique et réduction des
  répétitions. [À COMPLÉTER PAR L'AUTEUR : chapitres concernés.] La
  problématique, les hypothèses, le choix des protocoles et l'interprétation
  des résultats relèvent de l'auteur.

  \item \textbf{Figures et tableaux} : aide à l'écriture du script de
  génération des figures à partir des artefacts ; aucune valeur n'a été
  saisie manuellement dans une figure.

  \item \textbf{Bibliographie} : aucun outil n'a été utilisé comme source
  scientifique. Chaque référence de \texttt{references.bib} renvoie à une
  publication, une norme, une documentation ou un jeu de données consulté par
  l'auteur. [À COMPLÉTER PAR L'AUTEUR : confirmer, ou préciser si un outil a
  suggéré des références ensuite vérifiées.]

  \item \textbf{Résumés multilingues} : [À COMPLÉTER PAR L'AUTEUR : indiquer
  si l'Abstract et le Bokuse ont été traduits avec assistance, et qui les a
  relus.]
\end{itemize}

\section*{Mode de vérification}

\begin{itemize}
  \item Toute valeur numérique du mémoire est rattachée à un artefact
  identifié par un \texttt{run\_id} et un chemin du dépôt
  (tableau~\ref{tab:preuves-experimentales}) ; les valeurs proposées par un
  outil ont été comparées à ces artefacts avant intégration.
  \item Les descriptions du comportement du code ont été confrontées au code
  source et aux tests automatisés du dépôt.
  \item Les propositions de code ont été relues, exécutées et testées par
  l'auteur avant intégration. [À COMPLÉTER PAR L'AUTEUR : préciser la
  procédure, p.~ex. exécution de \texttt{pytest}.]
  \item Les rapports de révision assistée sont conservés dans le dossier
  \texttt{revision\_reports/} [À COMPLÉTER PAR L'AUTEUR : préciser s'ils sont
  publiés avec le dépôt ou conservés à titre privé].
\end{itemize}

\section*{Responsabilité}

L'auteur assume la responsabilité de la problématique, de la conception du
prototype, des protocoles, des résultats, de leur interprétation et du
contenu final du mémoire. Aucune responsabilité scientifique n'est transférée
à un système automatisé.
